\chapter{Generalized Method of Moments}
\label{cap:gmm}
% ── index ───────────────────────────────────────────────────────
\index{generalized method of moments|textbf}
\index{GMM|see{generalized method of moments}}
\index{gmm@\texttt{gmm()}|textbf}
\index{overidentification|textbf}
\index{J-test|textbf}
\index{moment conditions}

% ─────────────────────────────────────────────────────────────────────────────
\section{The framework}

The Generalized Method of Moments (GMM) unifies the estimators from previous
chapters under a single principle: identification via population \textbf{moment
conditions}.

A GMM model specifies $L$ moment conditions:
\[
  E\!\left[g_i(\beta)\right] = 0, \qquad g_i(\beta) \in \mathbb{R}^L
\]

OLS and IV are special cases:

\begin{center}
\begin{tabular}{lll}
\toprule
\textbf{Estimator} & \textbf{Moment condition} & \textbf{Situation} \\
\midrule
OLS   & $E[x_i\,\varepsilon_i] = 0$ & $L = K$, identified \\
IV/2SLS & $E[z_i\,\varepsilon_i] = 0$ & $L = K$, identified \\
GMM   & $E[z_i\,\varepsilon_i] = 0$ & $L > K$, overidentified \\
\bottomrule
\end{tabular}
\end{center}

With $L = K$ moments, $g(\hat\beta) = 0$ is solved exactly --- equivalent
to just-identified IV. With $L > K$, the system is overdetermined and no
$\beta$ zeroes all moments simultaneously; GMM minimizes a weighted
quadratic form:

\[
  \hat\beta_{\mathrm{GMM}}
  = \arg\min_\beta\; n\, \bar g(\beta)' \hat W^{-1} \bar g(\beta)
\]

where $\bar g(\beta) = \tfrac{1}{n}\sum_i g_i(\beta)$ and $\hat W$ is the
consistent estimate of the asymptotic variance of $\sqrt{n}\,\bar g(\beta)$.

% ─────────────────────────────────────────────────────────────────────────────
\section{The two-step estimator}

\begin{enumerate}
  \item \textbf{First step:} estimate $\hat\beta_1$ with $W = I$ (uniform weight).
  \item \textbf{Efficient weight matrix:} compute $\hat W = \tfrac{1}{n}
        \sum_i z_i z_i' \hat\varepsilon_i^2$ using first-step residuals.
  \item \textbf{Second step:} re-estimate $\hat\beta_2$ with $\hat W^{-1}$ ---
        the efficient GMM estimator.
\end{enumerate}

The two-step estimator is asymptotically efficient in the class of GMM
estimators with the same moment conditions.

% ─────────────────────────────────────────────────────────────────────────────
\section{J-test of overidentification}

When $L > K$, it is possible to test whether the excess moments are valid.
Hansen's statistic (\emph{J-test}):

\[
  J = n\, \bar g(\hat\beta)' \hat W^{-1} \bar g(\hat\beta)
  \;\xrightarrow{d}\; \chi^2(L - K)
  \quad \text{under } H_0: E[g_i(\beta_0)] = 0
\]

$H_0$ is the hypothesis that \emph{all} $L$ instruments are exogenous.
Rejecting $H_0$ indicates that at least one instrument is correlated with
$\varepsilon$ --- the model is misspecified or an instrument is invalid.

\begin{aviso}
  The J-test is necessary but not sufficient: failing to reject does not prove that all
  instruments are valid, only that the overidentifying restrictions are
  consistent with the data.
\end{aviso}

% ─────────────────────────────────────────────────────────────────────────────
\section{DGP Verification}

\subsection*{OLS, IV, and GMM side by side}

The DGP repeats the structure from Chapter~\ref{cap:iv}: $x$ is endogenous ($\mathrm{Cov}(x,\varepsilon)
\neq 0$), with two valid instruments $z_1$ and $z_2$. $\varepsilon$ is
orthogonalized against $z_1$ in-sample --- just-identified IV recovers
exactly; GMM with two instruments is overidentified ($L=2, K=1$).

\begin{lstlisting}[caption={GMM DGP --- biased OLS, exact IV, efficient GMM}]
seed(42)
input df
id
1
end
for i in 1..=10 { df = append(df, df) }
df |> filter(_n <= 1000)
generate df z1  = rnormal()
generate df z2  = rnormal()
generate df v   = rnormal()
let d1          = cov(df, z1, v) / variance(df, z1)
let d0          = mean(df, v) - d1 * mean(df, z1)
generate df eps = v - d0 - d1*z1
generate df x   = 0.8*z1 + 0.4*z2 + eps
generate df y   = 1.5 + 2.0*x + 0.7*eps
let m_ols = ols(y ~ x, df)
let m_iv  = iv(y ~ x, ~ z1, df)
let m_gmm = gmm(y ~ x, ~ z1 + z2, df)
display m_ols
display m_iv
display m_gmm
\end{lstlisting}

\begin{lstlisting}[language={}, basicstyle=\ttfamily\scriptsize,
  frame=none, numbers=none, backgroundcolor=\color{codbg},
  xleftmargin=0pt, xrightmargin=0pt, breaklines=false]
Dep. Variable: y  R²=0.9798  N=1000  F=48418.53
const   1.2757***  (0.0077)    x   2.3736***  (0.0108)

Dep. Variable: y  Estimator: 2SLS  N=1000
const   1.5000***  (0.0170)    x   2.0000***  (0.0264)

Dep. Variable: y  GMM Two-Step  N=1000  J=2.5012  Prob(J)=0.1138  DF=1
x0   1.4875***  (0.0145)    x1   2.0206***  (0.0226)
\end{lstlisting}

OLS is biased (endogeneity): $\hat\beta_1 = 2{,}37 \neq 2{,}0$. IV with $z_1$
recovers exactly $2{,}0000$ (in-sample orthogonalization). GMM with $z_1 + z_2$
converges to $2{,}0206$ --- consistent, not exact, because it uses two moments to
estimate one parameter. The J-test ($J = 2{,}50$, $p = 0{,}11$) does not reject
$H_0$: both instruments are valid.

\begin{nota}
  Hayashi displays the coefficients as \texttt{x0} (intercept) and \texttt{x1}
  (slope) in the GMM output. This is an internal display convention ---
  the coefficients correspond exactly to the order of variables in the formula.
\end{nota}

\subsection*{J-test rejects --- invalid instrument}

If $z_2$ is contaminated with $\varepsilon$, the J-test detects the violation:

\begin{lstlisting}[caption={Invalid instrument --- J-test rejects}]
// z2b correlated with eps: invalid instrument
generate df z2b = z2 + 0.5*eps
let m_bad = gmm(y ~ x, ~ z1 + z2b, df)
display m_bad
\end{lstlisting}

\begin{lstlisting}[language={}, basicstyle=\ttfamily\scriptsize,
  frame=none, numbers=none, backgroundcolor=\color{codbg},
  xleftmargin=0pt, xrightmargin=0pt, breaklines=false]
Dep. Variable: y  GMM Two-Step  N=1000  J=152.9417  Prob(J)=0.0000  DF=1
x0   1.3738***  (0.0089)    x1   2.2125***  (0.0124)
\end{lstlisting}

$J = 152{,}94$ ($p = 0{,}000$) rejects $H_0$ with high significance --- and the
coefficients are biased ($\hat\beta_1 = 2{,}21$). The J-test worked as an
alarm: before trusting the GMM estimates, check that $J$ does not reject.

% ─────────────────────────────────────────────────────────────────────────────
\section{Syntax}

\begin{lstlisting}
gmm(y ~ x1 + x2, ~ z1 + z2 + z3, df)
gmm(y ~ x1 + x2, ~ z1 + z2 + z3, df, cov=robust)
\end{lstlisting}

The first right-hand side (\lstinline{~ x1 + x2}) lists the regressors;
the second (\lstinline{~ z1 + z2 + z3}) lists the instruments. Exogenous
variables in $x$ must also appear as their own instruments.

\begin{lstlisting}[caption={Typical workflow with included exogenous variable}]
load "data.csv" as df

// educ is exogenous (instrument for itself)
// exp is endogenous; z1 and z2 are instruments for exp
let m = gmm(wage ~ educ + exp, ~ educ + z1 + z2, df)
display m
\end{lstlisting}

\section{Capturing the result}

\begin{lstlisting}
let m = gmm(y ~ x, ~ z1 + z2, df)
display m
let yhat = predict(m, "xb")
\end{lstlisting}
